\documentclass[12pt, a4paper]{article}
\usepackage[utf8]{inputenc}
\usepackage[T2A]{fontenc}
\usepackage{indentfirst, setspace}
\usepackage{tabularx, multirow}
\usepackage[normalem]{ulem}
\usepackage[style=russian]{csquotes}
\usepackage[english,russian]{babel}
\usepackage{hyperref}
\usepackage{ragged2e}
\usepackage{caption}
\usepackage{wrapfig}
\usepackage{amsmath}
\usepackage{tikz}
\makeatletter
\def\@biblabel#1{#1. }
\makeatother
\captionsetup{labelsep=endash}
\usepackage{listings}
\linespread{1.3}
\lstset{
  language=C++, % Язык программирования (можно выбрать другой)
  basicstyle=\ttfamily\small, % Основной стиль шрифта
  keywordstyle=\color{blue}, % Стиль ключевых слов
  breaklines=true,
  commentstyle=\color{green}, % Стиль комментариев
  stringstyle=\color{red}, % Стиль строк
  extendedchars=\true,
  showstringspaces=false,
}
\usepackage[left=20mm, top=15mm, right=15mm, bottom=20mm]{geometry}

\begin{document}
\begin{titlepage}
     \fontsize{12}{12}\selectfont

  {\centering

   \begin{bf}

    \begin{wrapfigure}{l}{19mm}
        
\includegraphics[width=22mm]{emblem.jpeg}
    \end{wrapfigure}


    \noindent Министерство науки и высшего образования Российской Федерации

    \noindent Федеральное государственное бюджетное образовательное учреждение высшего образования

    \noindent \enquote{Московский государственный технический унверситет

     \noindent имени Н.Э. Баумана

     \noindent (национальный исследовательский университет)}

    \noindent (МГТУ им. Н.Э. Баумана)

   \end{bf}
  }

  \vspace{0.4cm}

  {\setstretch{0.1}
   \noindent\rule{\textwidth}{1mm}
   \noindent\rule{\textwidth}{0.5mm}

  }

  \fontsize{14}{21}\selectfont

  \noindent\begin{tabularx}{\textwidth}{l >{\centering\arraybackslash}X}
   ФАКУЛЬТЕТ & \flqq Фундаментальные Науки\frqq \\ \cline{2-2}

   КАФЕДРА & ФН-12 \flqq Математическое моделирование\frqq \\ \cline{2-2}
  \end{tabularx}


  \vspace{1cm}


  \begin{center}
   \begin{bf}

    \fontsize{24}{36}\selectfont
    ОТЧЕТ

    \fontsize{20}{30}\selectfont
    ПО ЛАБОРАТОРНОЙ РАБОТЕ НА ТЕМУ:

    АВЛ дерево

   \end{bf}
  \end{center}

  \fontsize{14}{21}\selectfont
  \vspace{5cm}


  \noindent\begin{tabularx}{\textwidth}{ X >{\centering}p{4cm} p{1cm} c }
   Студент: & & & Подобедов В. В. \\ \cline{2-2} \cline{4-4}
   & \fontsize{10}{15}\selectfont дата, подпись & & \fontsize{10}{15}\selectfont Ф.И.О. \\
   Преподаватель: & & & Волкова Л. Л.\\ \cline{2-2} \cline{4-4}
   & \fontsize{10}{15}\selectfont дата, подпись & & \fontsize{10}{15}\selectfont Ф.И.О.
  \end{tabularx}

  \vspace{\fill}

  \begin{center}
   \it{Москва}, 2023
  \end{center}

  \thispagestyle{empty}
\end{titlepage}

\renewcommand{\contentsname}{Содержание}
\setcounter{page}{2}
\tableofcontents

\newpage
\begin{center}
   \section*{Введение} 
\addcontentsline{toc}{section}{Введение}
\end{center}
\justifying
\textbf{Цель лабораторной работы}: разработать программу для создания АВЛ-дерева.

Для достижения поставленной цели требуется решить следующие \textbf{задачи}.
\begin{enumerate}
\item Описать структуру АВЛ-дерева.
\item Описать алгоритмы добавления и удаления элемента, поиска элемента в дереве, балансировки дерева.
\item Разработать программу, предоставляющую пользователю выбор пункта меню, отображающую меню в цикле с постусловием, реализующую все описанные алгоритмы.
\item Выполнить тестирование программы.
\item Привести примеры работы программы.
\end{enumerate}

\newpage
\begin{center}
   \section{Аналитическая часть} 
\end{center}

\textbf{АВЛ-дерево}~\cite{book} --- это прежде всего двоичное дерево поиска, ключи которого удовлетворяют стандартному свойству: ключ любого узла дерева не меньше любого ключа в левом поддереве данного узла и не больше любого ключа в правом поддереве этого узла. Основная идея состоит в том, чтобы поддерживать высоты левого и правого поддеревьев каждого узла в балансе так, чтобы разница высот поддеревьев для каждого узла не превышала $1$ (то есть разница высот должна быть $-1$, $0$ или $1$). Каждый узел АВЛ-дерева содержит ключ, ссылки на левое и правое поддеревья, а также информацию о высоте. При вставке или удалении узлов дерево перебалансируется, чтобы сохранить свойство баланса высот. Это происходит путём поворотов узлов вокруг своих осей.

\newpage
\begin{center}
    \section{Конструкторская часть}
\end{center}
\subsection{Добавление элемента}
\begin{enumerate}
    \item Вставка нового элемента.
    \begin{itemize}
        \item[---] Создать функцию ``insert'', которая принимает корень дерева и новый ключ для вставки.
        \item[---] Если дерево пустое (корень --- ``nullptr''), создать новый узел с переданным ключом и сделать его корнем дерева.
        \item[---] Если ключ меньше ключа текущего узла, рекурсивно вставить его в левое поддерево.
        \item[---] Если ключ больше ключа текущего узла, рекурсивно вставить его в правое поддерево.
        \item[---] После вставки узла, обновить высоту текущего узла.
        \item[---] Вычислить балансировку ``balanceFactor'' для текущего узла (разница высот левого и правого поддерева).
        \item[---] Если балансировка узла нарушена (больше $1$ или меньше $-1$), выполнить балансировку путем вращений. 
    \end{itemize}
    \item Балансировка дерева.
    \begin{itemize}
        \item[---] Если баланс больше $1$ и вставка была в левое поддерево левого потомка, выполнить правый поворот для текущего узла.
        \item[---] Если баланс меньше $-1$ и вставка была в правое поддерево правого потомка, выполнить левый поворот для текущего узла.
        \item[---] Если баланс больше $1$ и вставка была в правое поддерево левого потомка, выполнить двойной поворот. Сначала левый поворот для левого потомка, затем правый поворот для текущего узла.
        \item[---] Если баланс меньше $-1$ и вставка была в левое поддерево правого потомка, выполнить двойной поворот. Сначала правый поворот для правого потомка, затем левый поворот для текущего узла.
    \end{itemize}
\end{enumerate}

\subsection{Удаление элемента}
\begin{enumerate}
    \item Начало удаления.
    \begin{itemize}
        \item[---] Функция ``remove'' принимает корень дерева ``root'' и ключ ``key'' для удаления.
        \item[---] Если корень дерева пуст, возвращается ``nullptr'', так как узла для удаления не существует.
    \end{itemize}
    \item Поиск узла для удаления.
    \begin{itemize}
        \item[---] Если ключ меньше ключа текущего узла, рекурсивно вызывается ``remove'' для левого поддерева.
        \item[---] Если ключ больше ключа текущего узла, рекурсивно вызывается ``remove'' для правого поддерева.
        \item[---] Если ключ совпадает с ключом текущего узла, происходит удаление узла.
    \end{itemize}
    \item Удаление узла.
    \begin{itemize}
        \item[---] Если у узла есть один потомок, узел заменяется этим потомком.
        \item[---] Если у узла два потомка, найденный узел с наименьшим ключом в правом поддереве заменяет удаляемый узел.
    \end{itemize}
    \item Балансировка после удаления.
    \begin{itemize}
        \item[---] После удаления узла, пересчитывается высота текущего узла. Проверяется баланс текущего узла.
        \item[---] Если баланс больше $1$ и левое поддерево текущего узла сбалансировано (balanceFactor(root->left) >= $0$), выполняется правый поворот.
        \item[---] Если баланс больше $1$ и левое поддерево текущего узла несбалансировано (balanceFactor(root->left) < $0$), выполняется левый поворот, а затем правый поворот.
        \item[---] Если баланс меньше $-1$ и правое поддерево текущего узла сбалансировано (balanceFactor(root->right) <= $0$), выполняется левый поворот.
        \item[---] Если баланс меньше $-1$ и правое поддерево текущего узла несбалансировано (balanceFactor(root->right) > $0$), выполняется правый поворот, а затем левый поворот. 
    \end{itemize}
\end{enumerate}

\subsection{Поиск элемента}
\begin{enumerate}
    \item Функция принимает указатель на корень дерева ``root'' и ключ ``key'' для поиска.
    \item Если корень дерева пуст, возвращается ``false'', так как узла с заданным ключом в дереве нет.
    \item Сравнение ключа.
    \begin{itemize}
        \item[---] Если ключ совпадает с ключом текущего узла, возвращается ``true'', так как узел с искомым ключом найден.
        \item[---] Если ключ меньше ключа текущего узла, рекурсивно вызывается ``search'' для левого поддерева.
        \item[---] Если ключ больше ключа текущего узла, рекурсивно вызывается ``search'' для правого поддерева.
    \end{itemize}
    \item Возврат результата.
    \begin{itemize}
        \item[---] Если узел с искомым ключом найден, возвращается ``true''.
        \item[---] Если узел не найден, возвращается ``false'', так как узла с искомым ключом в дереве нет. 
    \end{itemize}
\end{enumerate}

\newpage
\begin{center}
    \section{Технологическая часть}
\end{center}

Лабораторная работа была реализована на языке программирования C++. Кроме стандартной библиотеки <iostream>, которая предоставляет возможность пользоваься средствами ввода-вывода для стандартной консоли, была подключена библиотека <queue>. Она предоставляет реализацию очереди (queue) - структуры данных, основанной на принципе "первым пришёл — первым вышел" (FIFO - First-In-First-Out). Это означает, что элементы добавляются в конец очереди, а извлекаются из начала. Библиотека <algorithm> в C++ предоставляет множество стандартных функций для работы с контейнерами, итераторами и другими структурами данных.

Далее приведен листинг реализации трех алгоритмов (листинг~\ref{lst:label1}).
\begin{lstlisting}[caption={Исходный код программы для построения АВЛ-дерева},label={lst:label1}]
#include <iostream>
#include <algorithm>
#include <queue>

using namespace std;

class AVLTree {
private:
	struct Node {
		int key;
		Node* left;
		Node* right;
		int height;

		Node(int k) : key(k), left(nullptr), right(nullptr), height(1) {}
	};

	Node* root;

	int height(Node* node) {
		if (node == nullptr) return 0;
		return node->height;
	}

	int balanceFactor(Node* node) {
		if (node == nullptr) return 0;
		return height(node->left) - height(node->right);
	}

	Node* rightRotate(Node* root) {
		Node* x = root->left;
		Node* T2 = x->right;
		x->right = root;
		root->left = T2;

		root->height = max(height(root->left), height(root->right)) + 1;
		x->height = max(height(x->left), height(x->right)) + 1;

		return x;
	}

	Node* leftRotate(Node* root) {
		Node* y = root->right;
		Node* T2 = y->left;
		y->left = root;
		root->right = T2;

		root->height = max(height(root->left), height(root->right)) + 1;
		y->height = max(height(y->left), height(y->right)) + 1;

		return y;
	}

	Node* insert(Node* node, int key) {
		if (node == nullptr) return new Node(key);

		if (key < node->key) {
			node->left = insert(node->left, key);
		}
		else if (key > node->key) {
			node->right = insert(node->right, key);
		}
		else {
			return node;
		}
		
		node->height = 1 + max(height(node->left), height(node->right));

		int balance = balanceFactor(node);

		if (balance > 1 && key < node->left->key) {
			return rightRotate(node);
		}
		if (balance < -1 && key > node->right->key) {
			return leftRotate(node);
		}
		if (balance > 1 && key > node->left->key) {
			node->left = leftRotate(node->left);
			return rightRotate(node);
		}
		if (balance < -1 && key < node->right->key) {
			node->right = rightRotate(node->right);
			return leftRotate(node);
		}

		return node;
	}

	Node* minValueNode(Node* node) {
		Node* current = node;

		while (current->left != nullptr) {
			current = current->left;
		}

		return current;
	}

	Node* remove(Node* root, int key) {
		if (root == nullptr) return root;

		if (key < root->key) {
			root->left = remove(root->left, key);
		}
		else if (key > root->key) {
			root->right = remove(root->right, key);
		}
		else {
			if (root->left == nullptr || root->right == nullptr) {
				Node* temp = root->left ? root->left : root->right;

				if (temp == nullptr) {
					temp = root;
					root = nullptr;
				}
				else{ 
					*root = *temp; 
				}

				delete temp;
			}
			else {
				Node* temp = minValueNode(root->right);
				root->key = temp->key;
				root->right = remove(root->right, temp->key);
			}
		}

		if (root == nullptr) return root;

		root->height = 1 + max(height(root->left), height(root->right));

		int balance = balanceFactor(root);

		if (balance > 1 && balanceFactor(root->left) >= 0) {
			return rightRotate(root);
		}
		if (balance > 1 && balanceFactor(root->left) < 0) {
			root->left = leftRotate(root->left);
			return rightRotate(root);
		}
		if (balance < -1 && balanceFactor(root->right) <= 0) {
			return leftRotate(root);
		}
		if (balance < -1 && balanceFactor(root->right) > 0) {
			root->right = rightRotate(root->right);
			return leftRotate(root);
		}

		return root;
	}

	bool search(Node* root, int key) {
		if (root == nullptr) return false;

		if (key == root->key) {
			return true;
		}
		else if (key < root->key) {
			return search(root->left, key);
		}
		else {
			return search(root->right, key);
		}
	}

public:
	AVLTree() : root(nullptr) {}

	void insert(int key) {
		root = insert(root, key);
	}

	void remove(int key) {
		root = remove(root, key);
	}

	bool search(int key) {
		return search(root, key);
	}

	void printTreeLevelOrder() {
		if (root == nullptr) {
			cout << "Дерево пусто." << endl;
			return;
		}

		cout << "Дерево (по уровням):" << endl;
		queue<Node*> levelQueue;
		levelQueue.push(root);

		while (!levelQueue.empty()) {
			int levelSize = levelQueue.size();

			while (levelSize > 0) {
				Node* current = levelQueue.front();
				levelQueue.pop();

				cout << current->key << " ";

				if (current->left != nullptr) {
					levelQueue.push(current->left);
				}
				if (current->right != nullptr) {
					levelQueue.push(current->right);
				}

				levelSize--;
			}

			cout << endl;
		}
	}
};

int main() {
	setlocale(0, "");
	AVLTree tree;

	int choice, key;
	bool running = true;

	while (running) {
		cout << "Выберите операцию:" << endl;
		cout << "1. Добавить элемент" << endl;
		cout << "2. Удалить элемент" << endl;
		cout << "3. Поиск элемента" << endl;
		cout << "4. Вывести дерево (по уровням)" << endl;
		cout << "5. Выйти из программы" << endl;
		cout << "Введите выбор: ";
		cin >> choice;

		switch (choice) {
		case 1:
			cout << "Введите ключ для добавления: ";
			cin >> key;
			tree.insert(key);
			cout << endl;
			break;
		case 2:
			cout << "Введите ключ для удаления: ";
			cin >> key;
			tree.remove(key);
			cout << endl;
			break;
		case 3:
			cout << "Введите ключ для поиска: ";
			cin >> key;
			if (tree.search(key))
				cout << "Элемент найден." << endl;
			else
				cout << "Элемент не найден." << endl;
			cout << endl;
			break;
		case 4:
			tree.printTreeLevelOrder();
			break;
		case 5:
			running = false;
			break;
		default:
			cout << "Неверный выбор! Попробуйте еще раз." << endl;
		}
	}

	return 0;
}

\end{lstlisting}

\newpage
\begin{center}
    \section{Исследовательская часть}
\end{center}

Далее на рисунках представлен результат работы программы.

Результат работы программы, в котором создается АВЛ-дерево путем добавления туда элементов. Далее выводится созданное дерево, удаляется элемент и снова выводится АВЛ-дерево, приведен на рис.~\ref{img:grap1}. 
\begin{figure}[h!]
    \center{\includegraphics{test1.PNG}}
    \caption{Пример №1}
    \label{img:grap1}
\end{figure}

\newpage
\begin{figure}[h!]
    \center{\includegraphics{test1_1.PNG}}
    \caption{Пример №1 (продолжение)}
\end{figure}

\newpage
Создается АВЛ-дерево путем добавления туда элементов. Далее несуществующий элемент удаляется и ищется в АВЛ-дереве. Удаляется введенный элемент и выводится пустое дерево. Результат работы программы, приведен на рис.~\ref{img:grap2}. 
\begin{figure}[h!]
    \center{\includegraphics{test2.PNG}}
    \caption{Пример №2}
    \label{img:grap2}
\end{figure}

\newpage
Создается АВЛ-дерево путем добавления туда одинаковых элементов. Результат работы программы, приведен на рис.~\ref{img:grap3}. 
\begin{figure}[h!]
    \center{\includegraphics{test3.PNG}}
    \caption{Пример №3}
    \label{img:grap3}
\end{figure}

\newpage
Создается АВЛ-дерево путем добавления туда элементов. Далее существующий ищется в АВЛ-дереве и вывдится построенное АВЛ-дерево. Результат работы программы, приведен на рис.~\ref{img:grap4}. 
\begin{figure}[h!]
    \center{\includegraphics{test4.PNG}}
    \caption{Пример №4}
    \label{img:grap4}
\end{figure}

\newpage
\begin{center}
    \section*{Заключение}
    \addcontentsline{toc}{section}{Заключение}
\end{center}

Цель лабораторной работы достигнута. Разработана программа для создания АВЛ-дерева.

В результате выполнения лабораторной работы были выполнены все задачи.
\begin{enumerate}
\item Описана структура АВЛ-дерева.
\item Описаны алгоритмы добавления и удаления элемента, поиска элемента в дереве, балансировки дерева.
\item Разработана программа, предоставляющая пользователю выбор пункта меню, отображающая меню в цикле с постусловием, реализующая все описанные алгоритмы.
\item Выполнено тестирование программы.
\item Приведены примеры работы программы.
\end{enumerate}

\newpage
\begin{center}
\begin{thebibliography}{}
\addcontentsline{toc}{section}{Список используемых источников}
\bibitem{book} Никлаус Вирт. Алгоритмы и структуры данных. / Пер. с англ. Ткачев Ф. В. — М.: ДМК Пресс, 2010. — 272 с.
\end{thebibliography}
\end{center}

\end{document}
